% !TeX program = lualatex
% !TeX encoding = utf8
% !TeX spellcheck = uk_UA
% !TeX root = ../EMConspect.tex

%=========================================================
\chapter{Електричне поле статичної системи зарядів}\label{\currfilebase}
\graphicspath{{\currfilebase/Pictures}}
%=========================================================
%\epigraph{\Alexander  ...взаємне тяжіння електричної рідини, іменованої позитивною, і електричної рідини, іменованої зазвичай негативною,
%	полягає у зворотному відношенні квадратів відстаней...
%}{Шарль Огюстен Кулон}


%% --------------------------------------------------------
\section{Дипольний момент системи зарядів}
%% --------------------------------------------------------

Щоб описати розподіл зарядів у просторі та їхню взаємодію із зовнішнім електричним полем, часто недостатньо знати лише сумарний заряд системи. Важливу
роль відіграє просторове розташування зарядів, яке впливає як на створене ними поле, так і на їхню поведінку в зовнішньому полі. Для опису такої системи
вводять поняття \emphx{дипольного моменту}:

\begin{equation*}
	\vect{p} = \sum_{i=1}^{n}q_i\vect{r}_i,
\end{equation*}
де $\vect{r}_i$ --- радіус-вектор $i$-го заряду.


У загальному випадку ця
величина залежить від вибору початку відліку. Однак якщо система
загалом електронейтральна $\sum_{i=1}^{n} q_i = 0$, то
дипольний момент $\vect{p}$ системи визначений однозначно.


\emphx{Диполь} --- система, що складається з
двох точкових зарядів $q$, однакових за величиною і
протилежних за знаком.


\emphx{Плече диполя} --- вектор ($\vect{\ell}$), що йде від <<$-$>> заряду до
<<$+$>>, довжина якого дорівнює відстані між зарядами.


\begin{SCfigure}[][h!]
	\centering
	\localinput{elmomentum.tikz}
	\caption{Електричний диполь}
\end{SCfigure}
\emphx{Дипольним моментом} елементарного диполя називається вектор
\begin{equation}
	\vect{p} = q\vect{\ell}.
\end{equation}


\emphx{Жорсткий диполь} --- диполь із фіксованою відстанню між зарядами, що не
змінюється під зовнішніми впливами.


\emphx{Пружний диполь} --- диполь, де відстань між зарядами змінюється під зовнішніми
впливами, і дипольний момент може змінюватися.

Якщо необхідно визначити поле диполя на далеких відстанях, то зручно користуватись концепцією елементарного, або точкового диполя.
\emphx{Елементарним, або точковим диполем} називають граничну систему з $\ell \to 0$, $|q| \to \infty$ за кінцевої величини $p$.


%%% --------------------------------------------------------
%\subsection*{Дипольні моменти молекул}
%%% --------------------------------------------------------
%
%
%Дипольний момент виникає внаслідок різної електронегативності атомів, що складають
%молекулу, та
%їхнього розташування в просторі.
%
%
%Дебай (позначається як \textit{Д} або \textit{D}) --- позасистемна одиниця вимірювання
%електричного дипольного моменту молекул. Одиниця виміру названа на честь фізика П. Дебая.
%\begin{equation*}
%	1\ \text{Дебай} = 10^{-18}\ \text{Фр}\cdot\text{см}.
%\end{equation*}
%Більшість полярних молекул має дипольний момент порядку $1$ Д. Одиниця застосовується у
%фізичній хімії, атомній і молекулярній фізиці.
%
%\begin{table}[h!]
%	\centering
%	\caption{Дипольні моменти деяких полярних молекул}
%	\begin{tblr}{
%		colspec={X[c, m]X[c,m]},
%		row{1} = {c, m, fg=white, bg=themecolordark},
%		hlines,
%		vlines
%		}
%		Молекула                   & Дипольний момент (Дебай) \\
%		Вода (\ce{H2O})            & 1.84                     \\
%		Аміак (\ce{NH3})           & 1.46                     \\
%		Вуглекислий газ (\ce{CO2}) & 0.00                     \\
%		Хлороводень (\ce{HCl})     & 1.08                     \\
%		Метанол (\ce{CH3OH})       & 1.70                     \\
%	\end{tblr}
%\end{table}
%
%\begin{figure}[h!]
%	\centering
%
%	\localinput{moleculedipole.tikz}
%
%	\caption{Ілюстрація дипольних моментів деяких полярних молекул}
%\end{figure}
%
%\emphx{Неполярні молекули} --- це молекули, в яких електричні заряди розподілені рівномірно, і
%відсутній постійний дипольний момент. У таких молекулах атоми зазвичай мають однакову або
%близьку електронегативність, тому електрони діляться рівномірно між атомами.
%
%\begin{figure}[h!]\centering
%	\captionbox{Дипольний момент без зовнішнього поля $p = 0$.}[0.45\linewidth]{
%		\begin{tikzpicture}
%			\node[circle, draw=red, ball color=red!50, inner sep=0.1cm, text=red] (O) at (0,0)
%			{$+$};
%			\draw[ball color=blue!50, opacity=0.5] (O) circle (1);
%			\node[text=blue] at (135:0.6) {$-$};
%		\end{tikzpicture}
%	}
%	\captionbox{Дипольний момент виник під дією зовнішнього поля $p \neq 0$.}[0.45\linewidth]{
%		\begin{tikzpicture}[>=latex]
%			\node[circle, draw=red, ball color=red!50, inner sep=0.1cm, text=red] (O) at (0,0)
%			{$+$};
%			\draw[ball color=blue!50, opacity=0.5] ([xshift=-0.5cm]O) ellipse (1.5 and 0.75);
%			\node[text=blue] at (180:1) {$-$};
%
%			\foreach \y in {-1,-0.5,...,1}
%				{
%					\draw[->, red, opacity=0.3] (-3, \y) -- ++(5, 0);
%				}
%		\end{tikzpicture}
%	}
%\end{figure}
%
%
%Прикладом неполярних молекул є кисень (\ce{O2}), азот (\ce{N2}), метан (\ce{CH4}) та атоми
%інертних газів.
%
%Дипольний момент молекули має важливі практичні наслідки. Наявність сталого дипольного моменту визначає, чи здатна молекула поглинати електромагнітне випромінювання певних частот, обертаючись у зовнішньому полі. На цьому ґрунтується \emphx{мікрохвильова (обертальна) спектроскопія}: полярні молекули (наприклад, \ce{H2O}, \ce{CO}, \ce{HCl}) мають характерні обертальні спектри поглинання, тоді як неполярні молекули (\ce{O2}, \ce{N2}, \ce{CO2}) в обертальному русі електромагнітне випромінювання не поглинають. Завдяки цьому мікрохвильова спектроскопія дає змогу виявляти окремі молекули в міжзоряному середовищі та атмосферах планет, а інфрачервона спектроскопія --- ідентифікувати речовини за характерними коливально-обертальними смугами поглинання в хімічному аналізі.
%
%Полярність молекул, яку кількісно характеризує дипольний момент, визначає й розчинність речовин --- саме на цьому ґрунтується емпіричне правило <<подібне розчиняється в подібному>>: полярні розчинники (вода, спирти) добре розчиняють полярні речовини, тоді як неполярні розчинники (гексан, бензол) --- неполярні.



%% --------------------------------------------------------
\section{Поле точкового диполя}
%% --------------------------------------------------------


Знайдемо поле диполя на далеких відстанях $r$, тобто за умови $\ell \ll r$. У такому разі, ми можемо диполь вважати точковим. Для зручності розрахунків,
розташуємо
диполь на осі $Oy$, симетрично відносно початку координат (рис.~\ref{tikz:dipolefield}).

\begin{SCfigure}[0.75][h!]
	\centering
	\localinput{dipolefield.tikz}
	\caption{До виведення напруженості та потенціалу поля диполя}
    \label{tikz:dipolefield}
\end{SCfigure}

Поле точкового диполя можна знайти за допомогою принципу
суперпозиції
та з урахуванням нерівності $\ell \ll r$:
\begin{equation*}
	\Efield = \frac{(-q)\vect{r}_-}{r_-^3} +
	\frac{(+q)\vect{r}_+}{r_+^3}, \quad
	\begin{cases}
		\vect{r}_- =  \vect{r} + \vect{r}'_- \\
		\vect{r}_+ =  \vect{r} - \vect{r}'_+
	\end{cases}, \quad |\vect{r}'_-| = |\vect{r}'_+| = \frac{\ell}2.
\end{equation*}
\begin{equation}
	\tcbhighmath{
		\Efield = \frac{3(\vect{p}\cdot\vect{r})\vect{r}}{r^5} -
		\frac{\vect{p}}{r^3}.
	}
\end{equation}

Таким чином, напруженість поля точкового диполя зменшується з
відстанню за законом
\begin{equation*}
	E \sim 1/{r^3}.
\end{equation*}




%% --------------------------------------------------------
\section{Потенціал поля точкового диполя}
%% --------------------------------------------------------

%\begin{SCfigure}[2][h!]
%	\centering
%
%	\localinput{dipolefield.tikz}
%
%	\caption{До виведення потенціалу диполя.}
%\end{SCfigure}

Потенціал поля точкового диполя можна знайти за допомогою принципу
суперпозиції та з урахуванням нерівності $\ell \ll r$:
\begin{equation*}
	\phi = \frac{(-q)}{r_-} +
	\frac{(+q)}{r_+}, \quad
	\begin{cases}
		\vect{r}_- =  \vect{r} + \vect{r}'_- \\
		\vect{r}_+ =  \vect{r} - \vect{r}'_+
	\end{cases}, \quad |\vect{r}'_-| = |\vect{r}'_+| = \frac{\ell}2.
\end{equation*}
\begin{equation}
	\tcbhighmath{
		\phi = \frac{\vect{p}\cdot\vect{r}}{r^3}.
	}
\end{equation}

Таким чином, потенціал поля точкового диполя зменшується з
відстанню за законом
\begin{equation*}
	\phi \sim 1/{r^2}.
\end{equation*}



%% --------------------------------------------------------
\section{Потенціал системи зарядів на далеких відстанях}
%% --------------------------------------------------------
Розглянемо систему зарядів $\{q_i\}$, що займає певну скінченну область у просторі (рис.~\ref{tikz:far_potential}).

%---------------------------------------------------------
\begin{figure}[h!]\centering
    \localinput{far_potential.tikz}
    \caption{До визначення потенціалу на далеких відстанях}
    \label{tikz:far_potential}
\end{figure}
%---------------------------------------------------------


Потенціал поля в точці спостереження $P$ дорівнює сумі потенціалів, створюваних усіма зарядами:
\begin{equation*}
	\phi(\vect{r}) = \sum_i \frac{q_i}{|\vect{r} - \vect{r}_i|},
\end{equation*}
де $\vect{r}_i$ --- радіус-вектор $q_i$-го заряду, а
$\vect{r}$ --- радіус-вектор точки спостереження. Початок системи координат вибираємо всередині системи зарядів і знайдемо потенціал на відстанях, що
набагато більші за розмір області простору, в якій локалізуються заряди ($r_i \ll r$). Тоді в першому наближенні:
\begin{equation*}
    \frac{1}{|\vect{r} - \vect{r}_i|} \approx \frac1r \left(1 + \frac{\vect{r}\ \vect{r}_i}{r^2}\right).
\end{equation*}



Отже, $\phi(\vect{r})$ можна наближено записати як:
\begin{equation*}
	\phi(\vect{r}) \approx \sum_i \frac{q_i}r \left(1 + \frac{\vect{r}\ \vect{r}_i}{r^2}\right) = \frac{q}r + \frac{ \vect{p}\cdot\vect{r}}{r^2}.
\end{equation*}


Тут уведено повний заряд $q$ й дипольний момент $\vect{p}$ системи:
\begin{equation*}
    q = \sum_i {q_i}, \quad \vect{p} = \sum_i q_i\vect{r}_i.
\end{equation*}

Отже, потенціал поля системи зарядів в першому наближенні включає внесок від поля заряду (монополя) та поля диполя. Якщо система зарядів не
електронейтральна $q \neq 0$, то монопольна компонента значно більша за дипольний внесок. В іншому випадку дипольний момент може значно впливати на
потенціал, особливо у випадках, коли система має несиметричний розподіл зарядів.



%% --------------------------------------------------------
\section{Диполь в електричному полі}\label{sec:DipoleEnergy}
%% --------------------------------------------------------

%% --------------------------------------------------------
\subsection*{Момент сил, що діє на диполь в однорідному полі}
%% --------------------------------------------------------

\begin{SCfigure}[][h!]
	\centering
	\localinput{DipoleInUF.tikz}
	\caption{до виведення моменту сил, що діє на диполь в однорідному електричному полі}
    \label{tikz:Momentum_on_dipole_in_yniform_field}
\end{SCfigure}


Розглянемо диполь в однорідному електричному полі (рис.~\ref{tikz:Momentum_on_dipole_in_yniform_field}). Результуюча сила, що діє на цей диполь дорівнює
нулю: $\vect{F} = (+q)\Efield + (-q)\Efield = 0$.
Однак цей диполь може орієнтуватися у такому полі, оскільки на нього може діяти обертовий момент сил. Розрахуємо його, використовуючи означення,
відоме з курсу механіки:
\begin{equation*}
    \vect{M}_O = \vect{r}_+ \times \vect{F}_+ +   \vect{r}_- \times \vect{F}_- = \vect{r}_+ \times (+q)\Efield +   \vect{r}_- \times (-q)\Efield =
    q(\vect{r}_+ - \vect{r}_-)\times \Efield.
\end{equation*}
Оскільки $q(\vect{r}_+ - \vect{r}_-) = \vect{p}$, то остаточно маємо:


\begin{equation*}
	\tcbhighmath{
		\vect{M} = \left[\vect{p}\times\Efield\right].
	}
\end{equation*}


Бачимо, що момент сил (пари сил) не залежить від вибору точки $O$. Цей момент орієнтує диполь в однорідному полі.


%% --------------------------------------------------------
\subsection*{Сила, що діє на диполь у неоднорідному полі}
%% --------------------------------------------------------

Розрахуємо силу, що діє на диполь у неоднорідному полі (рис.~\ref{tikz:dipole_in_nonuniform_field}).

\begin{equation*}
	\vect{F} = (+q)\Efield_+ + (-q)\Efield_- = q(\Efield_+ - \Efield_-) =
	q\Delta\Efield.
\end{equation*}

\begin{SCfigure}[1.5][h!]
	\centering
	\localinput{dipinnonuniformfield.tikz}
	\caption{Диполь у неоднорідному полі}
    \label{tikz:dipole_in_nonuniform_field}
\end{SCfigure}

Для точкового диполя різницю $\Delta\Efield$ можна записати використовуючи формулу \eqref{eq:vec_diff} у вигляді:
\begin{equation*}
	\Delta\Efield = \Efield(\vect{r} + \vect{\ell}) - \Efield(\vect{r})= \ell_x
	\frac{\partial\Efield}{\partial x} +
	\ell_y \frac{\partial\Efield}{\partial y} + \ell_z
	\frac{\partial\Efield}{\partial z}.
\end{equation*}
\begin{equation}\tcbhighmath{
		\vect{F} = p_x \frac{\partial\Efield}{\partial x} +
		p_y \frac{\partial\Efield}{\partial y} + p_z \frac{\partial\Efield}{\partial z}
		=
		(\vect{p}\cdot\vect{\nabla})\Efield.
	}
\end{equation}

Якщо \emphx{диполь перебуває в неоднорідному полі} --- виникає \emphx{результуюча сила}, що діє
на диполь як  ціле, що \emphx{втягуватиме диполь в область сильного поля}.

% --------------------------------------------------------
\subsection*{Енергія жорсткого диполя в електричному полі}
%% --------------------------------------------------------

Енергія точкового заряду $q$ у зовнішньому полі дорівнює $U = q\phi$, де $\phi$ ---
потенціал поля в точці знаходження заряду $q$. Диполь --- це система з двох зарядів, тому
його енергія в зовнішньому полі $U= q\phi_+ + (-q)\phi_- = q (\phi_+ - \phi_-)$. Для різниці можна використати формулу~\eqref{eq:sc_diff}:
\begin{equation*}
	U= q\phi_+ + (-q)\phi_- = q (\phi_+ - \phi_-) = q \vect{\ell}\cdot\vect{\nabla}\phi.
\end{equation*}
Оскільки $-\vect{\nabla}\phi = \Efield$, а $\vect{p} = q\vect{\ell}$:
\begin{equation*}
	\tcbhighmath{
		U = -\vect{p}\cdot\Efield.
	}
\end{equation*}


Мінімальну енергію диполь має в положенні $\vect{p} \uparrow\uparrow \Efield$ (положення стійкої рівноваги). У разі відхилення з цього положення виникає
момент зовнішніх сил, що повертає диполь до положення рівноваги.
\begin{figure}[h!]
	\centering

	\localinput{dipoleenergy.tikz}

	\caption{Положення диполя в однорідному електричному полі. Ліворуч (\subref{tikz:parallel}) --- положення стійкої рівноваги з мінімальною енергією.
		Праворуч (\subref{tikz:angle}) --- нестійке
		положення.}
\end{figure}



%% --------------------------------------------------------
\subsection*{Енергія пружного диполя в електричному полі}
%% --------------------------------------------------------

Під дією електричного поля плече диполя збільшується на таку величину $\ell$, що сила з боку
поля врівноважується силою пружності
$
	k\ell = qE
$.
Звідси знаходимо, що дипольний момент $p = q\ell$ залежить від прикладеного поля за законом:
\begin{equation*}
	p = \frac{q^2E}{k} = \alpha E
\end{equation*}
Коефіцієнт $\alpha$ називається \emphx{поляризованістю молекули}.


Потенціальна енергія деформації дорівнює
$
	U = \frac{k\ell^2}2
$.
Оскільки $k = \frac{qE}{\ell}$, одержимо
$
	U = \frac{k\ell^2}2 = \frac{q\ell E}2 = \frac{pE}2
$.
Тут вектор дипольного моменту, індукований зовнішнім полем, паралельний електричному полю
$\vect{p}\uparrow\uparrow \Efield$. Тому останню формулу можна переписати у векторній формі:

\begin{equation}
	\tcbhighmath{
		U = \frac12\ \vect{p}\cdot\Efield.
	}
\end{equation}